\section{Decorations and restriction of cycles}\label{sec:chains}

The frame construction produces cycles whose terms end at projective
sign groups. The applications require two further operations. First,
we adjoin commuting elementary generators and cancel the additional
boundary faces. Second, we pass from a group of automorphisms to a
specified subgroup while retaining a nonzero full-flag coefficient.
Both operations concern actual coefficients of simplices. In
particular, parameters describing different apartments are not counted
as independent unless their images in the chain group are independent.

The use of coefficient equations to cancel faces is shared with the
constructible and admissible cycle methods of
\citet{DiazRamos2018,DiazRamosMazza2022}. Here we state the required
injectivity separately and verify it in each group application; the
decoration counts alone do not provide it.

\subsection{Split chains and their boundary}

Write $\widetilde C_*(P)$ for the augmented rational chain complex of
a finite poset $P$, with increasing order as the orientation of every
simplex. Let $S$ and $D$ be commuting elementary abelian $p$-subgroups
of a finite group $M$, and suppose $S\cap D=1$. Given flags
\[
 u=[U_1<\cdots<U_a],\qquad
 v=[V_1<\cdots<V_b]
\]
in $S$ and $D$, respectively, put $U_0=V_0=1$. An $(a,b)$-shuffle
is a path $\gamma$ from $(0,0)$ to $(a,b)$ whose steps are $(1,0)$
or $(0,1)$. Write $(i_k,j_k)$ for its vertex after $k$ steps, and
let $\nu(\gamma)$ count pairs consisting of a $D$-step followed
later by an $S$-step. Define
\begin{equation}\label{eq:chain-split-shuffle}
 u\boxtimes v
 =\sum_{\gamma}(-1)^{\nu(\gamma)}
 [U_{i_1}V_{j_1}<\cdots<U_{i_{a+b}}V_{j_{a+b}}].
\end{equation}
Each displayed inclusion is strict, because the two factors have
trivial intersection. The origin is omitted. Empty flags are allowed:
if $a=0$ or $b=0$, there is one path and this construction is the
identity on the other flag. Extend \eqref{eq:chain-split-shuffle}
bilinearly to chains. A chain of degree $a-1$ and a chain of degree
$b-1$ give a chain of degree $a+b-1$.

This is the product shuffle of Lemma~\ref{lem:product-shuffle},
specialized to elementary groups. That lemma gives
\begin{equation}\label{eq:chain-shuffle-boundary}
 \partial(u\boxtimes v)
  =(\partial u)\boxtimes v+(-1)^a u\boxtimes(\partial v).
\end{equation}
Here and below the empty simplex has degree $-1$ and zero boundary.
Iterating the construction gives the
sum over paths with any fixed number of kinds of steps. Its sign
counts inversions between the ordered kinds, so the iteration is
associative.

Suppose that $D=D_1\times\cdots\times D_t$, where the $D_i$ are
labelled subgroups of order $p$. Its coordinate apartment, coned
at $D$, is the chain
\begin{equation}\label{eq:chain-decoration-apartment}
 a(D_1,\ldots,D_t)=
 \sum_{\pi\in\mathfrak S_t}\sgn(\pi)
 [D_{\pi(1)}<D_{\pi(1)}D_{\pi(2)}<\cdots<D].
\end{equation}
Set $a(\varnothing)=[\,]$. Interior faces cancel in pairs,
and the remaining faces give the exact formula
\begin{equation}\label{eq:chain-decoration-boundary}
 \partial a(D_1,\ldots,D_t)
 =\sum_{i=1}^t(-1)^{i-1}
    a(D_1,\ldots,\widehat D_i,\ldots,D_t).
\end{equation}
For example, when $t=2$ the boundary is $[D_2]-[D_1]$.
If $x$ is a coned apartment in $S$, then $x\boxtimes a(D_1,\ldots,D_t)$
is the apartment for the union of its basis with the decoration
basis, up to the fixed choice of orientation of $x$. More generally,
\eqref{eq:chain-split-shuffle} defines this operation on the
\emph{coefficients} of any chain $x$, without choosing an expression
of $x$ as a sum of apartments.

\subsection{The coefficient spaces to be counted}

Fix an integer $h\geq2$, and put
\[
 r=h-1,\qquad q_h=(h-1)!.
\]
A frame $\mathcal F$ with $h$ summands has a projective sign group
$S_{\mathcal F}$ of rank $r$. Denote by $V_{\mathcal F}$ the space
of actual coned frame-chain coefficients at $S_{\mathcal F}$ from
Section~\ref{sec:frames}. Thus $V_{\mathcal F}$ is a subspace of the
chain group on full flags ending at $S_{\mathcal F}$, every member
has no boundary face retaining $S_{\mathcal F}$, and
\begin{equation}\label{eq:chain-local-dimension}
 \dim V_{\mathcal F}\leq q_h.
\end{equation}
In the abstract frame model equality holds. The inequality is the
only local bound needed here.

The precise output we use from Theorem~\ref{thm:frame-bound} is the
following. For the set $\mathfrak F$ of all allowed frames of one
of the spaces specified there, it provides a subspace
\begin{equation}\label{eq:chain-frame-input}
 Z\ \subseteq\ \bigoplus_{\mathcal F\in\mathfrak F}V_{\mathcal F},
 \qquad
 \dim Z\geq F\,|\mathfrak F|q_h,
 \qquad
 \sum_{\mathcal F\in\mathfrak F}\partial x_{\mathcal F}=0
       \quad(x\in Z).
\end{equation}
The last equality uses the coarsening identifications of that
theorem. Whenever the frame model is realized by elementary
subgroups, we require these identifications to be valid subgroup
identifications. A disjoint union of compatible spaces is also
allowed, provided the same lower bound $F$ holds for every space.
The direct sum of their spaces $Z$ then satisfies the same estimate.
For example, the all-type estimates with field parameter at least
five permit $F=(17/25)^h$; for symmetric $h=3$ they permit $F=3/8$.

Equation~\eqref{eq:chain-frame-input} counts the image on coned
flag coefficients, not the dimension of the word functions that
produce it. It still retains the frame labels in the direct sum.
If different full data can give the same subgroup flag after
decorations are adjoined, their possible coincidence must be
checked separately. The next lemma states this last requirement
as an injectivity hypothesis.

\begin{lemma}[Ordinary decorations]\label{lem:decorations}
Let $M$ be a finite group, let $p$ be a prime, and let $h\geq2$ and
$t\geq0$. Consider a finite nonempty set $\mathcal X$ of data
\[
 (\mathcal F,\delta),\qquad
 \delta=(D_1,\ldots,D_t),
\]
where $S_{\mathcal F}\leq M$ is a projective sign group of rank
$r=h-1$, the labelled $D_i\leq M$ have order $p$, and
\[
 E_{\mathcal F,\delta}
   =S_{\mathcal F}\times D_1\times\cdots\times D_t
\]
is an internal direct product. Suppose the following hold.
\begin{enumerate}
\item The local coefficient space $V_{\mathcal F}$ satisfies
\eqref{eq:chain-local-dimension} and is the same whenever the same
frame occurs with different decorations. For each fixed tuple
$\delta$, its compatible frames form a set $\mathfrak F_\delta$
with a subspace $Z_\delta$ satisfying
\eqref{eq:chain-frame-input}, with a common real number $F>0$.
The boundary equality remains valid in the subgroup poset of $M$.
\item Put $\mathcal Z=\bigoplus_\delta Z_\delta$. The linear map
\begin{equation}\label{eq:chain-detection-map}
 \Phi:\mathcal Z\longrightarrow
       \widetilde C_{r+t-1}(\Ap(M)),\qquad
 \Phi(x)=\sum_{(\mathcal F,\delta)\in\mathcal X}
 x_{\mathcal F,\delta}\boxtimes a(\delta)
\end{equation}
is injective.
\item For every $i\in\{1,\ldots,t\}$, each partial datum
$(\mathcal F,D_1,\ldots,\widehat D_i,\ldots,D_t)$ that occurs
has at least $k_i>0$ completions in $\mathcal X$.
\end{enumerate}
Write $N=|\mathcal X|$. Then the image of $\Phi$ contains a space
of cycles of dimension at least
\begin{equation}\label{eq:chain-decoration-dimension}
 Nq_h\left(F-\sum_{i=1}^t\frac1{k_i}\right).
\end{equation}
In particular, if $F>\sum_i1/k_i$, there is a nonzero cycle of
degree $r+t-1$, supported on full flags ending at the rank-$r+t$
groups $E_{\mathcal F,\delta}$. At least one such full flag has
a nonzero actual coefficient. When $t=0$, the sum is zero and
there are no completion conditions.
\end{lemma}

\begin{proof}
The fixed-tuple spaces give
\[
 \dim\mathcal Z
 \geq Fq_h\sum_\delta|\mathfrak F_\delta|=FNq_h.
\]
Use \eqref{eq:chain-shuffle-boundary} on
\eqref{eq:chain-detection-map}. For fixed $\delta$, the terms
$(\partial x_{\mathcal F,\delta})\boxtimes a(\delta)$ sum to
zero by \eqref{eq:chain-frame-input}. This use of that equality is
legitimate: multiplication by subgroups of the fixed decoration
group is a well-defined operation on each common coarsening face.

It remains to cancel the terms obtained from
\eqref{eq:chain-decoration-boundary}. For each partial datum
\[
 y=(\mathcal F,D_1,\ldots,\widehat D_i,\ldots,D_t)
\]
impose the vector equation
\begin{equation}\label{eq:chain-partial-equation}
 \sum_{\substack{(\mathcal F,\delta)\in\mathcal X\\
                  \delta\text{ completes }y}}
           x_{\mathcal F,\delta}=0
       \quad\hbox{in }V_{\mathcal F}.
\end{equation}
The boundary sign $(-1)^{r+i-1}$ is common to this entire sum.
The remaining decoration apartment is also common, so
\eqref{eq:chain-partial-equation} cancels its contribution.
These equations are sufficient even if different partial data
have further coincident faces.

The following dimension count is the vector-valued analogue of the
scalar constructible-cycle count in
\citet[equation~(17), arXiv version~1]{DiazRamos2018}.
Each full datum has exactly one partial datum of kind $i$.
Since each such partial datum has at least $k_i$ completions,
there are at most $N/k_i$ of them. Each vector equation has
rank at most $q_h$ by \eqref{eq:chain-local-dimension}. Their
common kernel in $\mathcal Z$ consequently has dimension at least
\eqref{eq:chain-decoration-dimension}. Its image consists of
cycles, and injectivity of $\Phi$ preserves this dimension.
Every summand in \eqref{eq:chain-detection-map} is a full flag
of the asserted rank, so a nonzero image has the stated coefficient.
\end{proof}

Here is a useful way to check the injectivity assumption in
Lemma~\ref{lem:decorations}. Suppose a full group $E$ determines
$S_{\mathcal F}$ and its frame $\mathcal F$. Suppose also that
$D=D_1\cdots D_t$ determines the labelled axes, for instance because
they lift a fixed ordered basis of a quotient on which $D$ maps
isomorphically. Inspect a split full flag which first grows through
$D_1<D_1D_2<\cdots<D$ and then through
\[
 DU_1<\cdots<DU_r=E
\]
for a sign flag $[U_1<\cdots<U_r=S_{\mathcal F}]$.
If another datum contributes to this flag, its full group gives
the same sign group and frame. At the vertex $D$, its split form
is $U V$ with $U\leq S_{\mathcal F}$ and $V$ contained in its
decoration complement. Since $D\cap S_{\mathcal F}=1$, necessarily
$U=1$. The rank of $D$ then forces the whole alternative complement
to equal $D$. The labelled axes agree by hypothesis. The coefficient
of the inspected flag is therefore the original sign coefficient,
up to the nonzero sign of the unique shuffle taking all decoration
steps first. These flags detect every coordinate in the direct sum.
The applications will verify the recovery of $S_{\mathcal F}$,
$\mathcal F$, and the labelled axes in their particular groups.

\subsection{A rank-two decoration with three lines}

For the symplectic construction the additional elementary group is
$T\cong C_2^2$. Its three order-two subgroups are not supplied
with a preferred pair of coordinate axes. Instead we use its full
two-dimensional coefficient space
\begin{equation}\label{eq:chain-quaternion-space}
 Q_T=
 \left\{\sum_{\ell<T,\ |\ell|=2}c_\ell[\ell<T]:
             \sum_\ell c_\ell=0\right\}.
\end{equation}
For $q=\sum c_\ell[\ell<T]\in Q_T$, direct calculation gives
\begin{equation}\label{eq:chain-quaternion-boundary}
 \partial q=-\sum_\ell c_\ell[\ell].
\end{equation}
The two-dimensional space in \eqref{eq:chain-quaternion-space}
is spanned by differences of its three flags, equivalently by the
apartments from pairs of independent lines. The name in the next
lemma refers to its later realization by a projective quaternion
group; the chain statement needs only $T\cong C_2^2$.

\begin{lemma}[Quaternion decoration]\label{lem:quaternion-decoration}
Let $M$ be a finite group, $h\geq2$, $r=h-1$, and
$\varepsilon\in\{0,1\}$. Consider a finite nonempty set
$\mathcal X$ of data $(\mathcal F,T)$ if $\varepsilon=0$, or
$(\mathcal F,T,B)$ if $\varepsilon=1$, where
\[
 E=S_{\mathcal F}\times T\times B,
 \qquad \rk S_{\mathcal F}=r,\quad T\cong C_2^2,
 \quad \rk B=\varepsilon.
\]
For $\varepsilon=0$, set $B=1$ and omit it from the data.
For each fixed $(T,B)$, assume its compatible frames have a space
$Z_{T,B}$ satisfying \eqref{eq:chain-frame-input}, with a uniform
lower bound $F>0$ and local bound $q_h=(h-1)!$. Assume that the map
\begin{equation}\label{eq:chain-quaternion-detection}
 \bigoplus_{(T,B)} Z_{T,B}\otimes Q_T
 \longrightarrow \widetilde C_{r+1+\varepsilon}
                       (\mathcal A_2(M))
\end{equation}
which sends a simple tensor to the sum of the split chains
$x_{\mathcal F,T,B}\boxtimes q\boxtimes[B]$ is injective.
When $B=1$, replace $[B]$ here by the empty simplex.

For every occurring partial datum consisting of a frame, a
retained line $\ell<T$, and $B$ when present, assume there are
at least $k_T>0$ completions to full data. If $\varepsilon=1$,
assume also that every occurring $(\mathcal F,T)$ has at least
$k_B>0$ completions to a datum $(\mathcal F,T,B)$.
Then the image in \eqref{eq:chain-quaternion-detection} contains
a cycle space of dimension at least
\begin{equation}\label{eq:chain-quaternion-dimension}
 2|\mathcal X|q_h
 \left(F-\frac{3}{2k_T}-\frac{\varepsilon}{k_B}\right),
\end{equation}
where the last term is omitted if $\varepsilon=0$.
If the expression in parentheses is positive, there is a
nonzero cycle on full flags ending at groups of rank
$r+2+\varepsilon=h+1+\varepsilon$.
\end{lemma}

\begin{proof}
Put $N=|\mathcal X|$. The domain of
\eqref{eq:chain-quaternion-detection} has dimension at least
$2FNq_h$. Its sign-direction boundaries cancel for each fixed
$(T,B)$ and each coefficient in $Q_T$, by
\eqref{eq:chain-frame-input} and
\eqref{eq:chain-shuffle-boundary}.

For a retained-line partial datum $(\mathcal F,\ell,B)$, take
the coefficient $c_\ell$ in \eqref{eq:chain-quaternion-space}
from each of its completions and sum the resulting vectors in
$V_{\mathcal F}$. Impose that this sum be zero. This is a
linear condition of rank at most $q_h$. By
\eqref{eq:chain-quaternion-boundary}, these conditions cancel
all boundary terms in which the $T$-part is reduced to a line.
There are exactly $3N$ incidences between full data and their
retained-line partial data. Thus at most $3N/k_T$ such partial
data occur, costing at most $3Nq_h/k_T$ dimensions.

If $B$ is present, its boundary is the empty simplex. For each
partial datum $(\mathcal F,T)$ impose zero sum of the full
vectors in $V_{\mathcal F}\otimes Q_T$ over its $B$-completions.
Each such condition has rank at most $2q_h$, and there are at
most $N/k_B$ such data. This costs at most $2Nq_h/k_B$ further
dimensions. The common kernel therefore has dimension at least
\eqref{eq:chain-quaternion-dimension}. Its image consists of
cycles, and the assumed injectivity preserves its dimension.
The split construction shows the asserted ranks and flag lengths.
\end{proof}

The use of two parameters at $T$ is essential to this count.
There are three retained-line boundary conditions but only one
sign coefficient vector at each such partial datum. Dividing
the resulting loss by the initial factor two gives
$3/(2k_T)$. To verify the detection hypothesis, the applications
inspect split flags through $T$ and each of its three lines;
when $B$ is present they begin at $B$ and then pass through
$B\ell<BT$. These flags recover the coefficients in $Q_T$
as well as the sign-chain coefficients, once the full data
are recovered from the corresponding subgroups.

\subsection{Opposite chambers and a nonzero restricted coefficient}

We now prove the restriction statement. Its input is a homogeneous
cycle on full flags of one fixed rank. This condition will let us
take a local building cycle at a specified top group. The following
elementary building fact supplies a complement on a supported flag.

Its conclusion has the classical apartment-basis interpretation of
\citet[Theorem~1]{Solomon1969} and
\citet[Proposition~4.5 and equation~(4.6)]{Bjorner1984}: coefficients
on opposite chambers detect a top cycle. We retain a direct panel
argument, including the augmented low-rank cases, before proving the
restriction statement that uses the resulting complement.

\begin{lemma}[A supported opposite chamber]\label{lem:opposite-chamber}
Let $V$ be an $m$-dimensional vector space over $\F_p$, with
$m\geq1$, and let $\mathcal B(V)$ be its poset of proper nonzero
subspaces. Let
\[
 0=F_0<F_1<\cdots<F_m=V,\qquad \dim F_j=j,
\]
be a fixed full flag with its endpoints adjoined. Every nonzero
cycle $z\in\widetilde C_{m-2}(\mathcal B(V))$ has a supported
chamber $W_1<\cdots<W_{m-1}$ such that, on putting $W_0=0$
and $W_m=V$,
\begin{equation}\label{eq:chain-opposition}
 \dim(W_i\cap F_j)=\max\{0,i+j-m\}
       \quad(0\leq i,j\leq m).
\end{equation}
In particular, if $0<c<m$ and $F_{m-c}=K$, then
$W_c$ is a complement to $K$ in $V$.
\end{lemma}

\begin{proof}
When $m=1$, the building is empty and its unique augmented
chamber is the empty simplex; the assertion holds directly.
Suppose $m\geq2$. Write $a_W$ for the coefficient of a chamber
$W=(W_1<\cdots<W_{m-1})$ in $z$. If all its vertices except
the rank-$i$ vertex are fixed, the cycle equation is
\begin{equation}\label{eq:chain-panel}
 \sum_{W_{i-1}<X<W_{i+1},\ \dim X=i}
 a_{(W_1,\ldots,W_{i-1},X,W_{i+1},\ldots,W_{m-1})}=0.
\end{equation}
The common boundary sign $(-1)^{i-1}$ has been omitted.
For $m=2$ this is the augmentation equation at the empty
simplex. A supported completion of any such panel therefore
has another supported completion.

For a chamber $W$, set $r_{ij}=\dim(W_i\cap F_j)$. There is
a permutation $w\in\mathfrak S_m$ with
\begin{equation}\label{eq:chain-relative-permutation}
 r_{ij}=|\{k\leq i:w(k)\leq j\}|.
\end{equation}
Indeed, $r_{ij}-r_{i-1,j}$ is the dimension of the image of
$W_i\cap F_j$ in the one-dimensional space $W_i/W_{i-1}$.
As $j$ increases these images are nested, so the difference
changes from zero to one at a unique threshold $w(i)$.
Since $\sum_i(r_{ij}-r_{i-1,j})=j$, exactly $j$ thresholds
are at most $j$, proving that the thresholds form a permutation.

Choose a supported chamber maximizing the inversion number
of $w$. If $w$ is not the decreasing permutation, some
adjacent entries satisfy $a=w(i)<w(i+1)=b$. Keep its panel
fixed, and put $A=W_{i-1}$ and $C=W_{i+1}$. On the
two-dimensional quotient $C/A$ the fixed flag induces
\[
 H_j=\bigl(A+(C\cap F_j)\bigr)/A.
\]
Its dimension is $r_{i+1,j}-r_{i-1,j}$. Hence it is zero
for $j<a$, one fixed line $\ell$ for $a\leq j<b$, and
all of $C/A$ for $j\geq b$. For an intermediate subspace
$A<X<C$,
\[
 \dim(X\cap F_j)-\dim(A\cap F_j)
       =\dim\bigl((X/A)\cap H_j\bigr).
\]
Thus $X/A=\ell$ is the unique panel completion with the
two thresholds in the order $(a,b)$. Every other completion
has them in the order $(b,a)$, with all other thresholds
unchanged. Such a completion has one more inversion.
Equation~\eqref{eq:chain-panel} supplies another supported
completion, contradicting maximality. Therefore $w$ is
decreasing, and \eqref{eq:chain-relative-permutation} gives
\eqref{eq:chain-opposition}. The case $i=c$, $j=m-c$ gives
$W_c\cap K=0$, and their dimensions sum to $m$.
\end{proof}

\begin{lemma}[Restriction at a maximal index]\label{lem:restriction}
Let $M$ be a finite group, $H\leq M$, and $p$ a prime. Let
$m\geq1$ and let $\alpha\in\widetilde C_{m-1}(\Ap(M))$
be a cycle supported on full flags ending at rank-$m$
elementary abelian groups. Denote by $\mathcal E(\alpha)$
the set of top groups of flags with nonzero coefficient in
$\alpha$. Suppose that an integer $c$ satisfies
\[
 0\leq c<m,
 \qquad [E':E'\cap H]\leq p^c
      \quad(E'\in\mathcal E(\alpha)),
\]
and that a specified $E\in\mathcal E(\alpha)$ satisfies
$[E:E\cap H]=p^c$. Then there is a nonzero cycle
\[
 \beta\in\widetilde C_{m-c-1}(\Ap(H))
\]
supported on full flags of rank $m-c$, with a nonzero
coefficient at a full flag ending at $E\cap H$.
If $m-c=m_p(H)$, this cycle represents a nonzero homology class.
\end{lemma}

\begin{proof}
If $c=0$, every supported top group lies in $H$, so take
$\beta=\alpha$. Suppose henceforth $0<c<m$, and put
$K=E\cap H$.

Collect the terms of $\alpha$ ending at $E$ and delete their
last vertex. This gives a nonzero chain
$z_E\in\widetilde C_{m-2}(\mathcal B(E))$, identifying
$E$ with an $m$-dimensional $\F_p$-space. It is a cycle:
each face retaining the rank-$m$ vertex $E$ can receive a
contribution only from terms of $\alpha$ whose top is $E$.
Its coefficient equation in $\partial\alpha=0$ is exactly
the corresponding equation in $\partial z_E=0$. This includes
the augmentation equation when $m=2$.

Choose a full reference flag of $E$ through $K$, which has
dimension $m-c$. By Lemma~\ref{lem:opposite-chamber}, a full
flag with nonzero coefficient in $\alpha$ passes through a
rank-$c$ subgroup $R$ with
\begin{equation}\label{eq:chain-chosen-complement}
 E=R\times K.
\end{equation}
Fix the initial segment $\sigma=[R_1<\cdots<R_c=R]$ of
that flag, where $\rk R_i=i$.

Apply the saturated residue $R_\sigma$ of
Lemma~\ref{lem:saturated-residue}: retain the simplices whose
initial segment is $\sigma$ and delete those $c$ vertices,
with their original coefficients. The result is a chain in
the strict upper interval of $R$, and that lemma gives
\begin{equation}\label{eq:chain-restriction-residue}
 R_\sigma\partial=(-1)^c\partial R_\sigma.
\end{equation}
Its hypotheses hold because $\rk R_i=i$; no subgroup can
be inserted below $R_1$ or between consecutive vertices.
Consequently $R_\sigma\alpha$ is a cycle of degree $m-c-1$.

Every vertex $B$ in its support contains $R$ strictly and
lies in some supported top group $E'\in\mathcal E(\alpha)$.
Since $R\cap H=1$, the subgroup indices satisfy
\begin{equation}\label{eq:chain-index-sandwich}
 p^c=|R|\leq[B:B\cap H]
      \leq[E':E'\cap H]\leq p^c.
\end{equation}
All groups under consideration are subgroups of the
elementary abelian group $E'$, so these are ordinary
vector-space index inequalities. Equality throughout gives
\begin{equation}\label{eq:chain-fixed-complement-split}
 B=R\times(B\cap H).
\end{equation}
In particular $B\cap H\ne1$. Formula
\eqref{eq:chain-fixed-complement-split} holds for every
continuation vertex, with the \emph{same} fixed subgroup $R$.
It follows that
\[
 B\longmapsto B\cap H
\]
is injective on these vertices, with inverse $T\mapsto RT$
on its image. It preserves and reflects strict inclusions.

Apply this map to $R_\sigma\alpha$, and call the image
$\beta$. Since it is an injective simplicial map on the
subcomplex generated by the support, $\beta$ is a cycle and
no two distinct suffixes have been identified. The chosen
nonzero coefficient therefore survives, with its top vertex
changed from $E$ to $E\cap H$. Moreover a vertex of rank
$c+j$ becomes one of rank $j$ by
\eqref{eq:chain-fixed-complement-split}; hence every resulting
flag is full of rank $m-c$. This proves all chain assertions.
If that rank equals $m_p(H)$, there are no chains in the
next degree, so the nonzero cycle cannot be a boundary.
\end{proof}

The maximal-index hypothesis is used only in
\eqref{eq:chain-index-sandwich}. A constant index over all
supported top groups is therefore sufficient, but is not
required. Neither $m$ nor $m-c$ needs to be the maximum
elementary rank of its ambient group unless the final
top-degree conclusion is invoked. In the applications at
two, the nonzero saturated coefficient supplied by
Lemma~\ref{lem:restriction} is instead used in
Lemma~\ref{lem:propagation}.

Figure~\ref{fig:fixed-complement} illustrates the fixed-complement step
when $0<c<m$. Its middle and bottom rows are two descriptions of the
same suffix flag, which is why the intersection operation preserves
the chosen coefficient.
\begin{figure}[ht]
\centering
\begin{tikzpicture}[
  box/.style={draw,rounded corners=2pt,inner sep=7pt,
              text width=12.5cm,align=center,font=\small},
  transfer/.style={-{Stealth[length=2mm]},thick}]
\node[box] (full) at (0,0)
  {$[R_1<\cdots<R_c=R<RT_1<\cdots<RT_{m-c}=E]$};
\node[box] (suffix) at (0,-1.45)
  {$[RT_1<\cdots<RT_{m-c}=E]$};
\node[box] (image) at (0,-2.90)
  {$[T_1<\cdots<T_{m-c}=K]$};
\draw[transfer] (full)--node[right,font=\small,fill=white,inner sep=2pt]
  {residue at $[R_1<\cdots<R_c]$} (suffix);
\draw[transfer] (suffix)--node[right,font=\small,fill=white,inner sep=2pt]
  {intersection with $H$} (image);
\end{tikzpicture}
\caption{The chosen supported flag in Lemma~\ref{lem:restriction}, with
$E=R\times K$, $K=E\cap H$, and $0<c=\rk R<m$.
Here $T_1<\cdots<T_{m-c}=K$ is a full flag of subgroups of $K$.
The maximal-index hypothesis forces every continuation vertex $B$ in
the residue support to equal $R\times(B\cap H)$ with this same $R$.
Thus intersection with $H$ is injective on those suffix flags and
preserves the selected nonzero coefficient. The diagram depicts a
flag in that support, not an arbitrary flag in the ambient poset.
When $c=0$, the initial flag is empty, $R=1$, and the lemma uses the
original cycle in $H$.}
\label{fig:fixed-complement}
\end{figure}

